\chapter{Categorical Symmetry And Defect Fusion}
\label{ch:global-categorical-symmetry-defect-fusion}

\ConventionsMixed

This chapter formulates categorical symmetry through topological defects and
their fusion.  The object replacing a group is a category or higher category
of topological operators, and the action on local or extended observables is
defined by surrounding, ending, or fusing defects.

\section{Topological Defect Category}

Let \(\mathcal Q\) be a \(D\)-dimensional local QFT.  A codimension-one
topological defect \(D\) is an operator supported on a hypersurface such that
correlation functions are invariant under compactly supported deformations of
the hypersurface avoiding insertions.  Defects form a category
\(\mathcal D_{\mathcal Q}\) whose objects are boundary labels on small
transverse intervals and whose morphisms are junction operators.  Fusion is a
monoidal product
\[
  \otimes:\mathcal D_{\mathcal Q}\times\mathcal D_{\mathcal Q}
  \longrightarrow \mathcal D_{\mathcal Q}.
\]
If \(D_g\otimes D_h\simeq D_{gh}\) and every \(D_g\) is invertible, the data
reduce to an ordinary group symmetry.  Categorical symmetry keeps the same
fusion operation but allows noninvertible simple objects and nontrivial
junction spaces.

\section{Action On Local Operators}

Let \(D\) be a topological codimension-one defect in Euclidean space.  Place a
small sphere \(S^{D-1}\) carrying \(D\) around a local operator
\(\mathcal O(0)\).  Shrinking the sphere defines a linear map
\[
  \rho_D:\mathcal V_{\mathrm{loc}}\longrightarrow\mathcal V_{\mathrm{loc}},
\]
where \(\mathcal V_{\mathrm{loc}}\) is the vector space of local operators at
the point, completed in the topology appropriate to the QFT.  Fusion implies
\[
  \rho_{D_1\otimes D_2}=\rho_{D_1}\rho_{D_2}.
\]
When \(D\) is not invertible, \(\rho_D\) may be a projection, a sum of
endomorphisms, or a map with a nontrivial kernel.  Selection rules follow from
eigenvectors and images of these maps, not from characters of a group alone.

\section{Fusion Coefficients And Junctions}

Assume the defect category is semisimple with simple objects \(D_a\).  Fusion
has coefficients
\[
  D_a\otimes D_b\simeq \bigoplus_c N_{ab}^{\ \ c}D_c,
  \qquad N_{ab}^{\ \ c}\in\mathbb Z_{\ge0}.
\]
A choice of basis in the junction spaces
\[
  \operatorname{Hom}(D_a\otimes D_b,D_c)
\]
gives associator matrices \(F\).  The equality of the two ways of fusing four
defects is the pentagon equation.  This equation is obtained by comparing two
isotopic networks of topological junctions in a ball.  Reflection positivity
and unitarity impose a \(*\)-structure and positivity conditions on the
junction pairings.

\section{Gauging A Categorical Symmetry}

For a finite group symmetry, gauging sums over flat bundles.  For a finite
categorical symmetry, the analogous operation is a condensation of a
commutative separable algebra object
\[
  A\in\mathcal D_{\mathcal Q}.
\]
The algebra maps are
\[
  m:A\otimes A\to A,\qquad
  \eta:\mathbf 1\to A,
\]
and the Frobenius and separability identities encode the local moves of the
defect network.  The gauged theory has topological defects given by
\(A\)-modules or \(A\)-bimodules, depending on the codimension and the
placement.  This construction is a QFT operation only when the defect-network
sum or condensation is defined as a regulated insertion in correlation
functions.

\section{Finite Regular Algebra And Self-Dual Gauging}

The finite-group case is the minimal example in which the abstract word
``condensation'' becomes an explicit algebra object.  Let \(H\) be a finite
group acting by distinct invertible topological defects
\[
  \{D_h:h\in H\},\qquad D_g\otimes D_h\simeq D_{gh}.
\]
Assume for the moment that the associator phases in the pointed defect
subcategory have been trivialized.  A nontrivial \(3\)-cocycle associator is
an anomaly or discrete-torsion datum; it modifies the junction maps below and
must be carried explicitly.

\begin{definition}[Regular algebra of a finite symmetry]
\label{def:categorical-regular-algebra-finite-symmetry}
The regular algebra object of the untwisted pointed defect subcategory
generated by \(H\) is
\[
  A_H=\bigoplus_{h\in H}D_h .
\]
Its multiplication is the direct sum of the group-law junctions
\[
  m_{g,h}:D_g\otimes D_h\longrightarrow D_{gh},
\]
and its unit is the inclusion \(\eta_A:\mathbf 1=D_e\hookrightarrow A_H\).
Choose the normalized comultiplication
\[
  \Delta(D_k)=
  \frac1{|H|}
  \sum_{\substack{g,h\in H\\ gh=k}}
  D_g\otimes D_h ,
  \label{eq:regular-algebra-normalized-comultiplication}
\]
where the displayed formula means the corresponding sum of dual junction
maps in the chosen bases.
\end{definition}

\begin{proposition}[Regular algebra identities]
\label{prop:regular-algebra-frobenius-separable}
In the untwisted pointed defect category generated by \(H\),
\[
  A_H\otimes A_H\simeq |H|\,A_H,
  \qquad
  m\circ\Delta=\operatorname{id}_{A_H},
  \label{eq:regular-algebra-object-and-separability}
\]
and the Frobenius identity
\[
  \Delta\circ m
  =
  (m\otimes\operatorname{id})\circ
  (\operatorname{id}\otimes\Delta)
  =
  (\operatorname{id}\otimes m)\circ
  (\Delta\otimes\operatorname{id})
  \label{eq:regular-algebra-frobenius-identity}
\]
holds on every simple summand \(D_g\otimes D_h\).
\end{proposition}

\begin{proof}
The object identity follows from
\[
  A_H\otimes A_H
  \simeq
  \bigoplus_{g,h\in H}D_{gh}.
\]
For fixed \(k\in H\), the equation \(gh=k\) has exactly \(|H|\) solutions:
choose \(g\) freely and set \(h=g^{-1}k\).  Hence each \(D_k\) occurs
\(|H|\) times.

For separability, apply \(m\circ\Delta\) to \(D_k\).  The summands in
\eqref{eq:regular-algebra-normalized-comultiplication} all multiply back to
\(D_k\), and there are \(|H|\) of them, so the normalized coefficient gives
one copy of \(D_k\).

For the Frobenius identity, evaluate on \(D_g\otimes D_h\).  The left side is
\[
  \Delta(D_{gh})
  =
  \frac1{|H|}
  \sum_{\substack{c,d\in H\\ cd=gh}}D_c\otimes D_d .
\]
The middle expression gives
\[
  \frac1{|H|}
  \sum_{\substack{a,b\in H\\ ab=h}}D_{ga}\otimes D_b .
\]
The map \((a,b)\mapsto(c,d)=(ga,b)\) is a bijection from the set
\(\{(a,b):ab=h\}\) to the set \(\{(c,d):cd=gh\}\).  Thus the two sums agree.
The right expression is identical by the bijection
\((a,b)\mapsto(c,d)=(a,bh)\) from \(\{(a,b):ab=g\}\) to
\(\{(c,d):cd=gh\}\).
\end{proof}

\begin{proposition}[Self-dual finite gauging fusion ring]
\label{prop:self-dual-finite-gauging-fusion-ring}
Suppose anomaly-free \(H\)-gauging of \(\mathcal Q\) admits an equivalence
\(\Phi:\mathcal Q/H\to\mathcal Q\).  Let
\(I_H:\mathcal Q\to\mathcal Q/H\) be the gauging interface, and let
\(\mathcal N_H=\Phi\circ I_H\) be the resulting gauging defect.  Assume the
equivalence has been chosen so
that \(\mathcal N_H^\dagger\simeq\mathcal N_H\) and the \(H\)-defects are
absorbed on both sides of the gauging interface:
\[
  D_h\otimes\mathcal N_H\simeq
  \mathcal N_H\otimes D_h\simeq \mathcal N_H .
\]
Then the object-level fusion ring generated by \(D_h\) and \(\mathcal N_H\)
is
\[
  D_gD_h=D_{gh},\qquad
  D_h\mathcal N_H=\mathcal N_HD_h=\mathcal N_H,\qquad
  \mathcal N_H^2=\bigoplus_{h\in H}D_h .
  \label{eq:self-dual-finite-gauging-fusion-ring}
\]
It is associative.  The Frobenius--Perron dimensions are
\[
  d(D_h)=1,\qquad d(\mathcal N_H)=\sqrt{|H|}.
\]
For \(|H|>1\), \(\mathcal N_H\) is not invertible.
\end{proposition}

\begin{proof}
The first rule is the ordinary symmetry-defect fusion.  The second is the
absorption property of an \(H\)-defect at an \(H\)-gauging interface.  The
third is the self-adjoint version of the slab fusion
\[
  \mathcal N_H^\dagger\otimes\mathcal N_H\simeq A_H
\]
from the finite groupoid sum over the slab between the interface and its
reverse.

Associativity is now an object-level calculation.  If only \(D_h\)'s occur,
it is associativity of the group \(H\).  If exactly one \(\mathcal N_H\)
occurs, both bracketings reduce to \(\mathcal N_H\).  If two
\(\mathcal N_H\)'s occur, for example
\[
  (\mathcal N_H^2)D_g
  =
  \bigoplus_{h\in H}D_{hg}
  =
  \bigoplus_{k\in H}D_k
  =
  \mathcal N_H(\mathcal N_HD_g),
\]
where right multiplication by \(g\) permutes \(H\); the other placements are
the same.  If three \(\mathcal N_H\)'s occur, then
\[
  (\mathcal N_H^2)\mathcal N_H
  =
  \bigoplus_{h\in H}D_h\mathcal N_H
  =
  |H|\,\mathcal N_H
  =
  \mathcal N_H\bigl(\mathcal N_H^2\bigr).
\]

Frobenius--Perron dimension is additive on direct sums and multiplicative on
fusion products.  Thus \(d(D_h)=1\), and the last fusion rule gives
\[
  d(\mathcal N_H)^2=\sum_{h\in H}d(D_h)=|H|.
\]
The positive solution is \(\sqrt{|H|}\).  If \(\mathcal N_H\) were invertible,
its square would be a single simple invertible object rather than a direct sum
of \(|H|>1\) distinct simple defects.
\end{proof}

\begin{example}[Cyclic self-dual gauging fusion rings]
\label{ex:cyclic-self-dual-gauging-fusion-rings}
For \(H=\mathbb Z_N=\langle a:a^N=e\rangle\), the finite self-dual gauging
fusion ring has simple invertible objects \(D_{a^j}\) and one noninvertible
object \(\mathcal N_N\) obeying
\[
  D_{a^i}D_{a^j}=D_{a^{i+j}},
  \qquad
  D_{a^j}\mathcal N_N=\mathcal N_ND_{a^j}=\mathcal N_N,
  \qquad
  \mathcal N_N^2=\bigoplus_{j=0}^{N-1}D_{a^j}.
\]
The \(N=2\) case is the Ising Kramers--Wannier fusion rule.  For \(N>2\) the
same equations are only the fusion-ring part of a possible defect category:
associators, junction pairings, and the existence of a local QFT realizing
the required self-duality are additional data, not consequences of the fusion
ring alone.
\end{example}

\section{The Ising/Kramers--Wannier Defect}

The simplest noninvertible example is the defect category of the critical
two-dimensional Ising model.  We describe the categorical data abstractly,
because these data are what survive as the symmetry object.  Let
the simple codimension-one topological defects be
\[
  \mathbf 1,\qquad \eta,\qquad \mathcal N .
\]
Here \(\mathbf 1\) is the invisible defect, \(\eta\) is the invertible
\(\mathbb Z_2\) spin-flip defect, and \(\mathcal N\) is the
Kramers--Wannier duality defect.  Their fusion rules are
\begin{equation}
\begin{gathered}
  \eta\otimes\eta\simeq\mathbf 1,\qquad
  \eta\otimes\mathcal N\simeq
  \mathcal N\otimes\eta\simeq\mathcal N,\\
  \mathcal N\otimes\mathcal N\simeq\mathbf 1\oplus\eta .
\end{gathered}
  \label{eq:ising-kw-fusion-rules}
\end{equation}
The last line is the decisive one: \(\mathcal N\) has no inverse, because
fusing it with itself gives two simple summands rather than the identity.

\begin{proposition}[Quantum dimensions of the Ising defects]
\label{prop:ising-defect-quantum-dimensions}
The Frobenius--Perron dimensions of
\(\mathbf 1,\eta,\mathcal N\) are
\[
  d_{\mathbf 1}=1,\qquad d_\eta=1,\qquad d_{\mathcal N}=\sqrt2 .
\]
\end{proposition}

\begin{proof}
Frobenius--Perron dimension is the positive homomorphism from the fusion
semiring to \(\mathbb R_{>0}\).  Thus \(d_{\mathbf 1}=1\), and
\(\eta^2=\mathbf 1\) with positivity gives \(d_\eta=1\).  The last fusion
rule in \eqref{eq:ising-kw-fusion-rules} gives
\[
  d_{\mathcal N}^2=d_{\mathbf 1}+d_\eta=2,
\]
hence \(d_{\mathcal N}=\sqrt2\).
\end{proof}

The defect is therefore noninvertible in a quantitative sense: surrounding
an insertion by \(\mathcal N\) is not an automorphism of the local-operator
space.  At the critical point the diagonal rational CFT has three primary
families, which we denote
\[
  \mathbf 1,\qquad \varepsilon,\qquad \sigma .
\]
The topological line action on these primary sectors is diagonalized by the
Ising modular \(S\)-matrix
\begin{equation}
  S_{\rm Ising}
  =
  \frac12
  \begin{pmatrix}
    1&1&\sqrt2\\
    1&1&-\sqrt2\\
    \sqrt2&-\sqrt2&0
  \end{pmatrix}_{\mathbf 1,\varepsilon,\sigma}.
  \label{eq:ising-modular-s-matrix}
\end{equation}
Here the row label \(\varepsilon\) is the same invertible topological line
as the spin-flip defect \(\eta\), and the row label \(\sigma\) is the same
topological line as the Kramers--Wannier defect \(\mathcal N\).  We keep the
notations distinct because \(\varepsilon,\sigma\) are also conventional
names for primary sectors, whereas \(\eta,\mathcal N\) emphasize the defect
interpretation.
The line \(D_a\) acts on the primary sector \(i\) by the scalar
\begin{equation}
  \lambda_a(i)=\frac{S_{ai}}{S_{\mathbf 1 i}} .
  \label{eq:ising-line-eigenvalues}
\end{equation}
For \(a=\mathcal N\), this gives
\[
  \lambda_{\mathcal N}(\mathbf 1)=\sqrt2,\qquad
  \lambda_{\mathcal N}(\varepsilon)=-\sqrt2,\qquad
  \lambda_{\mathcal N}(\sigma)=0 .
\]
The zero on the spin sector is not a paradox: a noninvertible defect need not
act as a bijection on local fields.

\begin{proposition}[Fusion action on Ising local sectors]
\label{prop:ising-defect-action-fusion}
The eigenvalues \eqref{eq:ising-line-eigenvalues} furnish a representation
of the fusion algebra:
\[
  \lambda_a(i)\lambda_b(i)
  =
  \sum_c N_{ab}^{\ \ c}\lambda_c(i)
\]
for every defect label \(a,b\) and every primary sector \(i\).
\end{proposition}

\begin{proof}
For \(\eta\) the eigenvalues are
\[
  \lambda_\eta(\mathbf 1)=1,\qquad
  \lambda_\eta(\varepsilon)=1,\qquad
  \lambda_\eta(\sigma)=-1,
\]
so \(\lambda_\eta^2=1\), matching
\(\eta\otimes\eta\simeq\mathbf 1\).  Multiplication by \(\eta\) leaves the
\(\mathcal N\)-eigenvalues unchanged, because
\(\mathcal N\) has zero eigenvalue on \(\sigma\), where \(\eta\) differs from
\(\mathbf 1\).  Finally,
\[
  \lambda_{\mathcal N}(i)^2
  =
  \begin{cases}
  2,& i=\mathbf 1,\varepsilon,\\
  0,& i=\sigma,
  \end{cases}
  =
  1+\lambda_\eta(i),
\]
which is exactly the action of
\(\mathbf 1\oplus\eta\).  These identities are the three fusion rules
\eqref{eq:ising-kw-fusion-rules} evaluated on local sectors.
\end{proof}

\begin{remark}[Duality defect from gauging interface]
Let \(\mathcal N_{\mathbb Z_2}\) be the interface that gauges the spin-flip
\(\mathbb Z_2\) symmetry.  At the self-dual critical Ising point the gauged
theory is equivalent to the original theory after choosing the duality
identification.  Composing the gauging interface with this identification
gives the Kramers--Wannier defect \(\mathcal N\).  The general interface
identity of Chapter~\ref{ch:global-duality-defects-gauging-orbifold-data}
then gives
\[
  \mathcal N^\dagger\otimes\mathcal N
  \simeq
  \mathbf 1\oplus\eta ,
\]
which is the noninvertible fusion rule above.  Thus the noninvertibility is
not a slogan about duality; it is the algebraic memory of summing over
\(\mathbb Z_2\)-bundles in the gauging interface.
\end{remark}

\section{SymTFT Interpretation In The Rational Example}

For a rational two-dimensional CFT, the categorical symmetry can be packaged
by a three-dimensional topological theory whose bulk line category contains
the defect category and whose boundary condition is the CFT.  In the Ising
case the relevant bulk topological line category has simple lines
\(\mathbf 1,\eta,\mathcal N\) with the same fusion rules
\eqref{eq:ising-kw-fusion-rules}.  A topological defect in the CFT is
obtained by ending the corresponding bulk line on the boundary.  Fusion of
defects is then inherited from fusion of bulk lines before they end.

This statement is deliberately precise but limited.  It does not assert that
an arbitrary continuum QFT with a categorical symmetry has already been
constructed from a fully extended topological bulk and boundary condition.
Rather, in rational examples it supplies a concrete comparison map:
\[
  \text{bulk topological lines}
  \longrightarrow
  \text{boundary topological defects}
  \longrightarrow
  \operatorname{End}(\mathcal V_{\rm loc}).
\]
The first arrow is the boundary endpoint operation, and the second is the
surrounding action on local operators.  The Ising computation above verifies
the last map explicitly on the three primary sectors.

\section{Anomaly As Obstruction To Condensation}

An anomaly of a categorical symmetry is an obstruction to defining the defect
networks independently of choices of triangulation, ordering, framing, or
extension to one higher dimension.  Algebraically it appears as failure of
the pentagon, braiding, or trace identities in the proposed physical
category.  Geometrically it is encoded by an invertible theory in one higher
dimension whose boundary is the original QFT with its defect category.

If the anomaly is cancelled by coupling to an invertible bulk, a defect
network on the boundary is evaluated together with a bulk topological term.
If the anomaly is trivialized by a local counterterm, the defect category can
be modified by an equivalent associator or braiding datum.

\section{Beyond The Rational Example}

The Ising example is a controlled rational model, not a substitute for the
general interacting problem.  Chapter~\ref{ch:global-duality-defects-gauging-orbifold-data}
develops two further finite or algebraic examples: noninvertible defects
obtained from gauging a finite normal subgroup at a self-dual point, and
duality walls acting on electric-magnetic line-operator lattices.  These
examples still leave the genuinely hard continuum problem open: construct the
defect category, junction operators, fusion rules, and anomaly or
condensation data directly in interacting QFTs rather than importing them
from a topological or rational model.

The next chapter continues this development for finite one-form symmetries.
It writes the five-dimensional \(\mathbb Z_N\) symmetry TQFT as an explicit
cochain theory, and evaluates the finite higher-gauging condensation defect by
an actual cochain sum.  The point of separating the chapters is logical: this
chapter defines the categorical action and its simplest rational example,
while Chapter~\ref{ch:global-higher-group-symmetry-tft} constructs the
higher-gauging mechanism that produces four-dimensional noninvertible
condensation defects when the necessary continuum surface operators are
available.

\begin{openproblem}[Interacting categorical symmetries]
\label{op:interacting-categorical-symmetries}
Construct categorical-symmetry defect categories in interacting continuum QFTs
with local algebras, reflection-positive state spaces, fusion of extended
operators, and anomaly inflow, beyond finite topological models and rational
two-dimensional examples.
\end{openproblem}
